Although applications of plate simulation continue to expand, access remains largely limited to specialized, proprietary software—often prohibitively expensive for students or even smaller research groups. This work presents the theoretical background and initial implementation of MAGPIE, an open-source platform for simulating plate physics. MAGPIE provides intuitive visualization tools and a range of analysis methods, making it useful both as a teaching aid for instructors and as a research tool for academics and industry professionals. Its features include defining elastic boundary conditions along plate edges, visualizing and exporting eigenshapes and corresponding frequencies, running time-domain simulations of the plate equation for sound synthesis, and estimating the elastic constants of experimental plates. This final function or elastic constants forms the basis of a tool for soundboard analysis to aid luthiers and conservators of historical string and keyboard instruments. 


